\documentclass[11pt,a4paper]{article}

%---------------------------------------------------------------------
% PACKAGES
%---------------------------------------------------------------------
\usepackage[T1]{fontenc}
\usepackage[utf8]{inputenc}
\usepackage[a4paper,top=2cm,bottom=2cm,left=2cm,right=2cm]{geometry}
\usepackage{libertine}
\usepackage{microtype}
\usepackage{titlesec}
\usepackage{enumitem}
\usepackage{tabularx}
\usepackage{xcolor}
\usepackage{hyperref}

%---------------------------------------------------------------------
% LAYOUT
%---------------------------------------------------------------------
\setlength{\parindent}{0pt}
\renewcommand{\baselinestretch}{1.05}
\pagestyle{empty}

\definecolor{darkgray}{HTML}{2a2a2a}
\hypersetup{
    colorlinks=true,
    linkcolor=darkgray,
    urlcolor=darkgray,
    pdftitle={Marcel Kück - Curriculum Vitae},
    pdfauthor={Marcel Kück}
}

% Section heading style
\titleformat{\section}
    {\large\scshape\bfseries}
    {}{0pt}{}
    [\vspace{-0.35em}\noindent\rule{\linewidth}{0.5pt}]
\titlespacing*{\section}{0pt}{1.2em}{0.45em}

% Bullet list spacing
\setlist[itemize]{
    leftmargin=1.3em,
    itemsep=0.2em,
    topsep=0.3em,
    parsep=0pt,
    label=\textbullet
}

%---------------------------------------------------------------------
% CUSTOM COMMANDS
%---------------------------------------------------------------------
% \cventry{Organisation}{Location}{Dates}
\newcommand{\cventry}[3]{%
    \begin{tabularx}{\textwidth}{@{}Xr@{}}
        \textbf{#1}, \textit{#2} & #3 \\
    \end{tabularx}%
    \vspace{-0.35em}\par%
}

% \cvitem{Name}{Dates}  (no location)
\newcommand{\cvitem}[2]{%
    \begin{tabularx}{\textwidth}{@{}Xr@{}}
        \textbf{#1} & #2 \\
    \end{tabularx}%
    \vspace{-0.35em}\par%
}

%---------------------------------------------------------------------
% DOCUMENT
%---------------------------------------------------------------------
\begin{document}

%-- HEADER -----------------------------------------------------------
\begin{center}
    {\Huge\scshape\bfseries Marcel Kück}\\[0.5em]
    {\small
        Munich, Germany \enspace$\bullet$\enspace
        \href{mailto:kueck.marcel@gmail.com}{kueck.marcel@gmail.com} \enspace$\bullet$\enspace
        +49 172 273 8908\\[0.15em]
        \href{https://www.linkedin.com/in/marcel-kueck}{LinkedIn} \enspace$\bullet$\enspace
        \href{https://github.com/MarcelKueck}{GitHub} \enspace$\bullet$\enspace
        \href{https://www.marcelkueck.dev/}{marcelkueck.dev}
    }
\end{center}

%-- RESEARCH INTERESTS ----------------------------------------------
\section*{Research Interests}

Physical intelligence and embodied AI, with a particular focus on tendon-driven and musculoskeletal actuation. Interested in the intersection of engineered biological tissue and robotic control: how learned controllers can be coupled to soft, biologically inspired actuators. Prior work includes neural-network-based control of a tendon-driven humanoid arm, imitation learning for bimanual manipulation, and hardware--software integration for humanoid platforms.

%-- EDUCATION --------------------------------------------------------
\section*{Education}

\cventry{Technical University of Munich}{Munich, Germany}{10.2019 -- 05.2024}
B.Sc.\ Computer Science. Specialization in Robotics and Artificial Intelligence.
\begin{itemize}
    \item Bachelor's thesis: \textit{Implementation and Analysis of a Neural Network-Based Control Algorithm for Optimizing Joint-Space Trajectory Tracking for a Musculoskeletal Tendon-Driven Humanoid Robot Arm}. Conducted at Devanthro; supervised by the \href{https://www.ce.cit.tum.de/en/air/home/}{Chair of Robotics, Artificial Intelligence and Real-Time Systems} (Prof.\ Alois Knoll).
    \item Member of \href{https://www.tum-ai.com/}{TUM.AI}, Software Development department.
\end{itemize}

\vspace{0.4em}
\cventry{University of Oxford}{Oxford, United Kingdom}{05.2024 -- 07.2024}
Machine Learning Summer School. Coursework and research seminars on reinforcement learning, generative AI, and machine learning in medicine. Selective admission (7\%).

\vspace{0.4em}
\cvitem{Google Machine Learning Bootcamp Europe}{07.2022 -- 11.2022}
Three-month programme on applied machine learning. Selective admission ($\sim$10\%).

%-- RESEARCH EXPERIENCE ----------------------------------------------
\section*{Research Experience}

\cventry{Devanthro}{Munich, Germany}{03.2026 -- Present}
\textit{Affiliated Researcher (project-based)}
\begin{itemize}
    \item Conducting independent research on control and embodied AI for Robody 5.0, the next-generation humanoid platform built on the OpenArm bimanual robotic arm.
    \item Building an end-to-end pipeline from VR-based teleoperation through LeRobot policy training to on-hardware deployment.
    \item Investigating data-efficient imitation learning for bimanual manipulation.
\end{itemize}

\vspace{0.4em}
\cventry{Devanthro}{Munich, Germany}{10.2022 -- 08.2024}
\textit{Research Student \& Founder's Associate}
\begin{itemize}
    \item Conducted research on neural-network-based control for tendon-driven humanoid arms (Robody 3.0), culminating in the bachelor's thesis on joint-space trajectory tracking for a musculoskeletal tendon-driven arm.
    \item Co-developed Robody 4.0: assisted in mechanical assembly, contributed CAD work in Fusion 360, and authored a Python script for parametric design of the elastic hull (``skin'') of the Viper arm.
    \item Supported the founders across fundraising, strategy, and general operations as Founder's Associate.
\end{itemize}

%-- INDUSTRY EXPERIENCE ----------------------------------------------
\section*{Industry Experience}

\cventry{MDK Engineering}{Munich, Germany}{01.2026 -- Present}
\textit{Founder}
\begin{itemize}
    \item Solo AI and robotics consultancy providing AI integration and workflow automation as a service to SMB and enterprise clients across multiple sectors.
    \item Designs and delivers retrieval-augmented generation systems, conversational agents, and AI-powered internal tools tailored to client domains.
    \item Builds automation pipelines on n8n and Make to integrate LLM and ML capabilities into client business processes and operations.
\end{itemize}

\vspace{0.4em}
\cventry{Rechnungs-API \normalfont{\small(\href{https://www.rechnungs-api.de}{rechnungs-api.de})}}{Munich, Germany}{02.2026 -- Present}
\textit{Co-Founder \& CEO}
\begin{itemize}
    \item Co-founded a B2B invoicing API; built and operated with a two-person team.
    \item Active paying customers; approximately 1{,}000 daily visits.
    \item Responsible for product direction, technical architecture, and go-to-market.
\end{itemize}

\vspace{0.4em}
\cventry{ShareYourSpace \normalfont{\small(\href{https://www.shareyourspace.com}{shareyourspace.com})}}{Munich, Germany}{02.2025 -- 01.2026}
\textit{Co-Founder \& COO}
\begin{itemize}
    \item Co-founded a marketplace for shared and flexible office space, an Airbnb-style platform for the coworking sector.
    \item Assumed day-to-day operational and product leadership during the company's final year.
    \item Led the sale of the company and software platform to Pixida GmbH, an existing client.
\end{itemize}

\vspace{0.4em}
\cventry{Siemens AG}{Munich, Germany}{09.2021 -- 12.2023}
\textit{Software Engineer, \href{https://sgi.siemens.com/home}{Siemens SGI}}
\begin{itemize}
    \item Designed and shipped a recommender system serving 50{,}000+ digital resources to 300{,}000+ users.
    \item Built machine-learning and web features across the SGI platform.
    \item Stack: Python, NumPy, TensorFlow, SQL, TypeScript, JavaScript, Angular.
\end{itemize}

%-- PERSONAL PROJECTS ------------------------------------------------
\section*{Personal Projects}

\textbf{RustyRobots:} Open-source robotics framework for embodied AI in Rust; modular crates for perception, control, and simulation integration. \href{[link]}{[link]}

\vspace{0.25em}
\textbf{RustML:} Open-source machine-learning framework in Rust; tensor operations, automatic differentiation, and inference. \href{[link]}{[link]}

\vspace{0.25em}
\textbf{LeRobot SO-101 Manipulator:} End-to-end vision-language-action policy training and on-hardware deployment on the SO-101 arm using the LeRobot framework. \href{[link]}{[link]}

\vspace{0.25em}
\textbf{Roasterio:} Smart coffee-bean roaster combining a custom sensor stack with automatic roast-profile optimization.

\vspace{0.25em}
\textbf{LaMarcello:} Retrofit controller converting classic espresso machines into smart, app-connected appliances; PID temperature control and flow profiling.

\vspace{0.25em}
\textbf{Murph:} Online medical consultation platform pairing medical students with elderly users for accessible primary advice.

\vspace{0.25em}
\textbf{Marie Lou Coffee:} Personal coffee-roasting venture; small-batch artisanal roasting, exploring a commercial launch.

%-- TECHNICAL SKILLS -------------------------------------------------
\section*{Technical Skills}

\begin{tabularx}{\textwidth}{@{}lX@{}}
\textbf{Programming:} & Python, Rust, C++, TypeScript, JavaScript \\[0.15em]
\textbf{Machine Learning \& Robotics:} & PyTorch, TensorFlow, LeRobot, ROS\,2, NVIDIA Isaac Sim, NVIDIA Cosmos \\[0.15em]
\textbf{Hardware \& Tooling:} & Linux, Git, Docker, Fusion 360, SLA printing, Jetson, sensor integration \\
\end{tabularx}

%-- LANGUAGES --------------------------------------------------------
\section*{Languages}

German (native) \enspace$\bullet$\enspace English (fluent, C1--C2)

\end{document}